\documentclass[12pt]{article}
\usepackage[francais]{babel}
\usepackage[utf8]{inputenc}
\usepackage[T1]{fontenc}
\usepackage{a4wide}
\usepackage{eurosym}
\usepackage{comment}
\usepackage{tikz}
% - REMARK: do not use blanks or newlines between the libraries
\usetikzlibrary{arrows,decorations.pathmorphing,decorations.footprints,fadings,shadows,decorations.text,patterns,positioning,shapes,matrix,fit}
\usepackage{calc}
\usepackage{fp} %floating point
\usepackage{rotating}


\title{Gestion de projet\\ 
Cas \textit{Personal Fit}}
\author{Stéphane Genaud}
\begin{document}
\thispagestyle{empty}
\maketitle


\section{Présentation}
\paragraph{Le contexte} 

\noindent  4.5.6 est une enseigne de prêt à porter pour femme, dont la cible est
la clientèle 30-50 ans. La marque est présente sous la forme d'un réseau de
franchisés (95 boutiques en France). Le chiffre d'affaire moyen d'une boutique est de
\EUR{850k}, et le résultat net de \EUR{300k} (annuellement).

\paragraph{Organisation}

Sa direction est organisée comme suit:

\begin{figure}[!htb]
\begin{tikzpicture} [xscale=1,yscale=1]
	\node (pdg) at (-5,0) {\includegraphics[height=12mm] {../fig/user_executive.png}};
	\draw (-5,-.75) node[below]{{\footnotesize PDG}};
	\node (daf) at (-3,0) {\includegraphics[height=12mm] {../fig/exec_green_user.png}};
	\draw (-3,-1.25) node[below]{{\footnotesize DAF}};
	\node (dp) at (-0.5,0) {\includegraphics[height=12mm] {../fig/registered_user.png}};
	\draw (-0.5,-.75) node[below]{{\footnotesize Dir. Production}};
	\node (dp) at (2.5,0) {\includegraphics[height=12mm] {../fig/user_female2.png}};
	\draw (2.5,-1.25) node[below]{{\footnotesize Dir. Commerciale}};
	\node (codir) at (5,0) {\includegraphics[height=15mm] {../fig/user_group.png}};
	\draw (5,-.75) node[below]{{\footnotesize Comité de Direction}};
\end{tikzpicture}
\end{figure}

\paragraph{Résultat de l'évaluation et modifications du projet}
\noindent 
Il a été décidé de lancer le projet \textit{Personal Fit}. Après avoir sondé 
les franchisés, il ressort que 70 points de ventes pourraient s'équiper dans 
les trois mois à venir (septembre à décembre). Pour remporter l'adhésion des
franchisés, d'abord réticents, la direction a voté un budget spécifique de 
\EUR{1,5 M} qui sera distribué aux franchisés sous la forme d'un prêt à taux 
zéro remboursable sur 2 ans.
4.5.6 a également obtenu, grâce à une commande de 70 cabines, de 20\% de réduction
sur le prix initial de la prestation, à l'origine de \EUR{32k}.

Une étude de marché, dont les conclusions sont encourageantes, a également été
un argument fort pour remporter l'adhésion des franchisés. L'étude relève qu'une
majorité des sondés dans la catégorie d'âge 45-65 ans, c{\oe}ur de cible de la
marque, trouve l'offre d'un grand intérêt.

\begin{figure}
\begin{center}
\begin{tabular}{|c|c|c|c|c|}
\hline
intérêt	&	aucun  & peu  &  moyen 	& grand \\
\hline
15-24 ans	&	19\%	&  41\% &  25\%	& 15\% \\
25-34 ans	&	17\%	&  34\% &  30\%	& 19\% \\
35-44 ans	&	11\%	&  25\% &  39\%	& 21\% \\
45-64 ans	&	5\%	&  6\% &   35\%	& 54\% \\
65-84 ans	&	15\%	&  13\% &  42\%	& 30\% \\
\hline
\end{tabular}
\caption{Extrait de l'étude de marché: intérêt pour le concept par tranche d'âge}
\end{center}
\end{figure}




\paragraph{Lancement du projet}

La maîtrise d'ouvrage du projet à été confiée à la direction commerciale.
Le livrable du projet a été défini comme l'installation fonctionnelle des 
70 points de ventes. La date de fin prévue, comme dans l'idée initiale, 
reste le 31 décembre, tandis que que le début du projet est prévu dans 
quelques jours, le 1er août.
\begin{figure}[!htb]
\begin{center}
\begin{tikzpicture}
 	\node(chefprojet) at (-1,0) {\includegraphics[height=11mm] {../fig/user_female.png}};
	\node(dircommerciale) at (2,0) {\includegraphics[height=9mm] {../fig/user_female2.png}};
	\draw (-1,-.75) node[below]{{\footnotesize chef de projet}};
	\draw (2,-.75) node[below]{{\footnotesize dir. commerciale}};
	\draw[->,ultra thick] (chefprojet) -- (dircommerciale);
\end{tikzpicture}
\end{center}
\end{figure}





\section{Questions}

Le chef de projet doit maintenant établir l'ensemble des tâches à accomplir.
Etablissez 
\begin{enumerate}
\item une liste des prestations matérielles ou intellectuelles qu'il va
falloir prévoir pour réaliser le projet.
\item une liste des responsables ou prestataires externes qu'il va falloir 
interroger pour obtenir des informations complémentaires quant à ces 
prestations.
\item établissez une matrice comportant, en lignes, les prestations 
à réaliser, et en colonnes, les actions pour réaliser ces prestations. 
\end{enumerate}




%-------------------------------------------------------------------------------
%\begin{comment}

\section{Corrigé}

\paragraph{Matrice WBS}
~\\	


\begin{tabular}{|l||l|l|l|l|l|l|l|l|}
\hline
objet $\backslash$ action 	&\begin{sideways}livraison\end{sideways}   
					&\begin{sideways}maintenance\end{sideways}
					&\begin{sideways}installation \end{sideways}
					&\begin{sideways}conception \end{sideways}
					&\begin{sideways}fabrication\end{sideways}
					&\begin{sideways}adaptation\end{sideways}
					&\begin{sideways}audit\end{sideways}
					&\begin{sideways}formation\end{sideways}\\
\hline
cabines			&	x	   & x	   &	x		&	 	&		&		& 	&\\	
espace Personal Fit	&		   &  	   &  x		&		& 		& 		& x 	&\\
signalétique		&		   & 		   &  x		& x		& x		&		&	&\\			
chaine production		& 		   &		   &  		&		&		& x		& 	&\\	
chaine livraison		& 		   &		   &  		&		&		& x		& 	&\\	
tracabilité informatique& 		   &		   &  		&		&		& x		& 	&\\	
site internet		&		   & 		   &  x		& x		& x		&		&	&\\			
app. internet prise de mesure &	   &		   &  x		&		&		& 		&     &\\	
conseillers ventes	&		   &		   &			&		&		&		&	& x \\		
\hline
\end{tabular}

%\end{comment}
%---------------------------------------------------------------------
\end{document}
